\paragraph{\trContactTypeBoundary{}的情形。}\label{sub:contact-type-boundary-case}
接着\cref{sub:contact-type-boundary} 继续，我们讨论带\trContactTypeBoundary{}的%
\trExactSymplectic\trManifold{}的情形。给定这样一个\trManifold{} $(M,\varpi,\theta)$，可以把
$\mathscr{J}$ 的定义改为由一切在边界附近\trContact{}型的\trCompatible\trAlmostComplex{}%
结构所成的空间。这意味着，在所得的\trFloerDatum{}与\trPerturbationDatum{}概念中，
\trAlmostComplex{}结构族 $J(t)$ 或 $J(z)$ 在 $\partial M$ 附近不再取常值。由于
\cref{thm:family-convexity}，加上\trInhomogeneous{}项在边界附近为零这一事实，
方程\eqref{eq:floer-equation}、\eqref{eq:inhomogeneous-equation} 的解仍将留在
$M$ 的内部。以这种方式定义 $\scr{F}(M)^{\mathrm{pr}}$，就会发现它本质上不依赖于%
\trAlmostComplex{}结构的任何选取（其含义是：不同的选取当然构成一个\trCoherent{}系统）。
\begin{proposition}\label{thm:contact-type-embedding-quasi-iso}
设 $M$ 是带\trContactTypeBoundary{}的\trExactSymplectic\trManifold，$N$ 是带角的%
\trExactSymplectic\trManifold，而 $\iota\colon M\hookrightarrow\mathrm{int}(N)$ 是一个%
\trEmbedding，满足 $\iota^*\omega_N=\omega_M$、$\iota^*\theta_N=\theta_M+d(\text{函数})$。%
那么 $\scr{F}(M)^{\mathrm{pr}}$ 与 $\scr{F}(N)^{\mathrm{pr}}$ 中由位于 $\iota(\mathrm{int}(M))$
之内的那些\trExact\trLagrange\trSubmanifold{}组成的\trAinfty-\trFullSubcategory{}
\trQuasiisomorphic。
\end{proposition}
\begin{proof}
完全不触动 $\scr{F}(N)^{\mathrm{pr}}$，可以用一个在 $\partial N$ 附近为零的函数的%
\trDifferential{}来改变 $\theta_N$。这样，就把局面化归到了 $\iota^*\theta_N=\theta_M$ 的情形。
做完这一步之后，我们径直把 $\iota$ 看作一个包含。取定 $N$ 上一个\trAlmostComplex{}%
结构 $J$，它在 $\partial N$ 附近与 $I_N$ 相同，并且使 $J|M$ 在 $\partial M$ 附近%
\trContact{}型。在定义所论的 $\scr{F}(N)^{\mathrm{pr}}$ 的\trFullSubcategory{}时，
可以选取所有\trAlmostComplex{}结构在 $N\setminus\mathrm{int}(M)$ 的一个邻域内都等于 $J$，%
并取\trHamiltonian{}函数与\trPerturbationDatum{}使之在 $N\setminus\mathrm{int}(M)$
的一个邻域内为零。于是\cref{thm:relative-maximum-principle}适用，表明一切%
\trInhomogeneous\trPseudoholomorphic\trPolygon{}都留在 $\mathrm{int}(M)$ 之内。在这种情形，两个%
\trCategory{}实际上相等。对一般的数据选取，\trQuasiisomorphism{}则由前文关于\trWellDefinedness{}的%
诸考虑得出。
\end{proof}
\begin{corollary}\label{thm:cone-completion-quasiequivalence}
设 $M$、$N$ 是带\trContactTypeBoundary{}的\trExactSymplectic\trManifold，而
$\widehat{M}$、$\widehat{N}$ 是照\eqref{eq:cone-completion} 那样贴\trCone{}所得的%
非紧\trManifold。设存在一个\trExactSymplectic\trIsomorphism{}
$\hat{\phi}\colon\widehat{M}\to\widehat{N}$。那么诸\trCategory{}
$\scr{F}(M)^{\mathrm{pr}}$、$\scr{F}(N)^{\mathrm{pr}}$ 相互\trQuasiequivalence。
\end{corollary}
\begin{proof}
考察诸\trCompletion{}的有限部分。
\begin{equation*}
\widehat{M}^{\sigma}=M\cup_{\partial M}\bigl([0;\sigma]\times\partial M\bigr),
\qquad
\widehat{N}^{\tau}=N\cup_{\partial N}\bigl([0;\tau]\times\partial N\bigr),
\end{equation*}
取 $0\ll\sigma\ll\tau$ 使得 $N\subset\hat{\phi}(\widehat{M}^{\sigma})\subset
\widehat{N}^{\tau}$。首先，$M$ 的\trFukayaCategory{}\trQuasiequivalence{}于 $\widehat{M}^{\sigma}$ 的%
（同理，$N$ 的与 $\widehat{N}^{\tau}$ 的\trQuasiequivalence）。这一点容易看出：利用负%
\trLiouville{}\trFlow，可以把两个\trManifold{}互相同化，只差\trSymplectic{}形式及其%
1-形式原函数的一个\trConformal\trRescaling。若随即相应地\trRescaling\trFloerDatum{}与\trPerturbationDatum，%
两个\trFukayaCategory{}实际上就会相同。
现在应用\cref{thm:contact-type-embedding-quasi-iso}，它说 $\scr{F}(\widehat{M}^{\sigma})$
与 $\scr{F}(\widehat{N}^{\tau})$ 的某个\trFullSubcategory{}\trQuasiisomorphic。该\trSubcategory{}
显然包含位于 $N$ 自身之内的全部\trExact\trLagrange\trSubmanifold，而在\trIsotopy{}的意义下%
这就是全部，故我们便得到所要的\trQuasiequivalence。
\end{proof}
鉴于\cref{thm:contact-type-moser}，便有：
\begin{corollary}\label{thm:deformation-invariance-contact}
$\scr{F}(M)^{\mathrm{pr}}$ 在带\trContactTypeBoundary{}的\trExactSymplectic\trManifold{}%
类内 $M$ 的\trDeformation{}下不变。\noqed
\end{corollary}
\begin{remark}\label{thm:completion-contact-action-remark}
作为一个小的变体，也可以利用\trCompletion{}上自然的 \trAlmostComplex{}结构类
$\hat{J}$ 直接定义诸%
\trCategory{} $\scr{F}(\widehat{M})^{\mathrm{pr}}$。\cref{thm:completion-maximum-principle}
随即会表明 $\scr{F}(M)^{\mathrm{pr}}$ 与 $\scr{F}(\widehat{M})^{\mathrm{pr}}$
\trQuasiequivalence。这一做法的一个优点是：它给出在无穷远处\trContact{}型的%
\trExactSymplectic\trAutomorphism{}群 $\mathrm{Aut}(\widehat{M})$ 在
$\scr{F}(M)^{\mathrm{pr}}$ 上的一个\trCoherent{}作用。
\end{remark}
\begin{corollary}\label{thm:rounding-quasiequivalence}
设 $N$ 是带角的\trExactSymplectic\trManifold。照\cref{thm:rounding-lemma}那样%
\trRoundOffCorners，取 $\delta$ 充分小，得到带\trContactTypeBoundary{}的\trExactSymplectic\trManifold{}
$M\subset\mathrm{int}(N)$。那么 $\scr{F}(M)^{\mathrm{pr}}$ 与 $\scr{F}(N)^{\mathrm{pr}}$
\trQuasiequivalence。
\end{corollary}
这由\cref{thm:contact-type-embedding-quasi-iso}与每一\trExact\trLagrange\trSubmanifold{}%
都可移入 $M$ 之内这一事实立即得出。
\begin{corollary}\label{thm:deformation-invariance-corners}
在不计\trQuasiequivalence{}的意义下，带角\trExactSymplectic\trManifold{}的\trFukayaCategory{}%
在\trDeformation{}下不变。
\end{corollary}
\begin{proof}
取 $N$ 与一个\trDeformation{} $(\omega_N^s,\theta_N^s,I_N^s)$。细察
\cref{thm:rounding-lemma}的证明即可看出：存在 $M\subset\mathrm{int}(N)$，使得对一切 $s$，
$(\omega_N^s,\theta_N^s)$ 限制在 $M$ 上都给出带\trContactTypeBoundary{}的%
\trExactSymplectic\trManifold{}结构。\cref{thm:deformation-invariance-contact}表明
$\scr{F}(M)^{\mathrm{pr}}$ 不依赖于 $s$，而\cref{thm:rounding-quasiequivalence}则让人
把结果搬到 $\scr{F}(N)^{\mathrm{pr}}$ 上。
\end{proof}
\begin{corollary}\label{thm:isomorphism-invariance}
设 $\phi\colon N\to\widetilde{N}$ 是带角\trExactSymplectic\trManifold{}之间的一个%
\trIsomorphism。那么存在从 $\scr{F}(N)^{\mathrm{pr}}$ 到 $\scr{F}(\widetilde{N})^{\mathrm{pr}}$
的一个\trQuasiequivalence，它把每个\trExact\trLagrange\trSubmanifold{} $L$ 映到
$\phi(L)$。
\end{corollary}
在给出证明之前，应当先澄清一个可能的混淆之源。如果该\trIsomorphism{}与诸%
\trAlmostComplex{}结构\trCompatible——即在边界附近 $\phi_*(I_N)=I_{\widetilde{N}}$——%
那么\cref{sub:well-definedness} 中的形式论证就足以证明结果（事实上，它们给出一个%
稍强的陈述，因为所得的\trQuasiequivalence{}自动在相关\trFunctor\trCategory{}中%
\trIsomorphic{}的意义下是\trCanonical{}的）。一般情形的困难恰恰来自如下事实：我们在
\cref{sub:exact-symplectic-with-corners} 中定义的\trIsomorphism{}概念无视\trAlmostComplex{}结构。
\begin{proof}
\trRoundOffCorners，得到带\trContactTypeBoundary{}的\trExactSymplectic\trManifold{}
$M\subset\mathrm{int}(N)$。限制 $\phi|M$ 显然是一个\trExactSymplectic\trEmbedding。
先把\cref{thm:rounding-quasiequivalence}的\trQuasiequivalence{}求逆，再应用
\cref{thm:contact-type-embedding-quasi-iso}，便得到一个\trAFunctor
\begin{equation}\label{eq:isomorphism-quasi-equivalence}
\scr{F}(N)^{\mathrm{pr}}\longrightarrow\scr{F}(M)^{\mathrm{pr}}
\longrightarrow\scr{F}(\widetilde{N})^{\mathrm{pr}}
\end{equation}
它在\trCohomology{}上\trFullyFaithful。选取小的 $\delta$，便可安排 $\widetilde{N}$ 的任何%
\trExact\trLagrange\trSubmanifold{}经负\trLiouville\trFlow{}被移入 $\phi(\mathrm{int}(M))$
之内，于是\eqref{eq:isomorphism-quasi-equivalence} 是一个\trQuasiequivalence。按此处
的定义，它把 $L$ 映到 $\phi(\tilde{L})$，其中 $\tilde{L}$ 与 $L$ \trIsotopy。然而，
由于 $\phi(L)$ 与 $\phi(\tilde{L})$ 在 $H\scr{F}(\widetilde{N})^{\mathrm{pr}}$
中是\trIsomorphic{}的\trObject，一个纯代数的论证使我们可以修改该\trFunctor，%
使得 $L$ 真正被映到 $\phi(L)$。
\end{proof}
